\documentclass[11pt]{article}
\usepackage[margin=1.1in]{geometry}
\usepackage{amsmath}
\usepackage{amssymb}
\usepackage{amsthm}
\usepackage{booktabs}
\usepackage{longtable}
\usepackage[hidelinks]{hyperref}

\theoremstyle{definition}
\newtheorem{finding}{Finding}
\newtheorem*{verdictbox}{Verdict}

\newcommand{\z}{\zeta(5)}
\newcommand{\Ub}{\overline{U}}
\newcommand{\vG}{v_p^{\mathrm{G}}}
\newcommand{\gin}{\gamma_p^{\mathrm{in}}}
\newcommand{\gout}{\gamma_p^{\mathrm{out}}}
\newcommand{\Iout}{I_{\mathrm{out}}}

% verdict tokens for the summary table
\newcommand{\OK}{correct}
\newcommand{\IM}{imprecise}
\newcommand{\GP}{\textbf{gap}}
\newcommand{\FA}{\textbf{false}}
\newcommand{\NC}{not checked}

\tolerance=1200
\emergencystretch=2em

\title{\bf Referee audit of\\ ``$\zeta(5)$ is irrational'' (A.\ Fauzan, 17 September 2026)}
\author{Claude Opus 5 (1M context), Anthropic\\
\normalsize an AI system, run as an agentic workflow in Claude Code\\
\normalsize prepared for Dan Romik}
\date{22 September 2026}

\begin{document}
\maketitle

\section{Executive summary}

\subsection*{What the paper claims}

The preprint --- Aabir Fauzan, \emph{$\zeta(5)$ is irrational}, 17 September 2026, Zenodo,
\href{https://doi.org/10.5281/zenodo.22826418}{doi:10.5281/zenodo.22826418}; this audit is of version v1,
\href{https://doi.org/10.5281/zenodo.22826419}{doi:10.5281/zenodo.22826419}, and all page and equation
references are to that version (32 pp., single author) --- proves two things. \textbf{Theorem 1.1:} $\z$ is irrational.
\textbf{Corollary 1.2:} $|\z-a/b|>b^{-260}$ for all integers $a$ and all $b\ge b_0$, with
$b_0$ unspecified. Corollary 1.2 is not needed for Theorem 1.1.

\subsection*{How the proof is organised}

Fix $K=40n$, $N=3n$, $h=37n$, $\alpha=3/40$, $\lambda=37/40$, and a cutoff $M\ge 40$ with
$K\ge 200M^2$. Equations (2.2)--(2.3) define a $\mathbb{Q}$-linear functional
$\mu_X$ on rational functions with simple poles among the $-j^2$, with values in $\mathbb{Q}[X]$
and affine dependence on $X$. Proposition 2.2 shows that at $X=\z$ the functional
is integration against the positive weight $w(y)=\tfrac{(2\pi)^4}{12}y^5\sum_{\ell\ge1}\ell^4e^{-2\pi\ell y}$,
so the $h\times h$ Hankel matrix $G_K(X)=[\mu_X(D_N^6t^{i+j}/D_K)]$ is at $X=\z$ a positive
moment matrix; (2.9) computes the leading coefficient of $\Delta_K=\det G_K$ by a Vandermonde
identity and shows $\deg\Delta_K=h$ exactly. Sections 3--5 clear denominators: a $p$-adic
extension of $\mu_X$ to a Tate algebra (Lemma 3.1), a distribution relation (Lemma 3.2) and a
crude small-prime bound (Lemma 3.3) yield local determinant valuations $\gin$ (Proposition 4.1,
for $K/M<p\le K/3$) and $\gout$ (Proposition 4.3, for $p>K/3$); (5.1)--(5.3) assemble these into
an explicit positive rational $m_{K,M}$ with $Q_{K,M}=m_{K,M}S_K\Delta_K\in\mathbb{Z}[X]$
(Proposition 5.1), and the prime number theorem converts the local bounds into
$\limsup K^{-2}\log m_{K,M}\le A_M$ (Proposition 5.2). Section 6 bounds the real side:
Andr\'eief's identity (6.10) writes $\Delta_K(\z)$ as an $h$-fold integral with a squared
Vandermonde, which after the scaling $y_i=K\sqrt{t_i}$ is a discrete logarithmic energy in the
external field $V$ of (6.1); comparison with an explicit sum of sixteen arcsine measures
(Lemma 6.1) together with negativity of zero-mass logarithmic energy (Lemma 6.2) gives
$\log F_K(\z)\le\Ub K^2+24K\log K+200K$ (Proposition 6.3). Section 7 adds the two constants:
$A_{200}+\Ub<-139/8000<0$, which is (2.7); then $b^{37n}Q_n(a/b)$ is a positive integer tending
to zero, and Theorem 1.1 follows.

\subsection*{The verdict}

\begin{verdictbox}
\textbf{We found no error.} Nine independent audit tracks, each followed by a separate
adversarial verification pass, checked every numbered statement of the paper that reduces to a
computation --- most of them in exact rational or $p$-adic arithmetic, several at 25--6000
significant digits. Not one of the paper's statements came out false. The claimed identities
that can be checked exactly are exact; the claimed inequalities that can be checked hold,
several of them with equality; and the two decisive constants reproduce digit for digit from
the paper's own data.

\textbf{The paper as written is not a complete proof.} It has one gap: on p.~10 the
inequality $b_a\le 6\alpha x+3$ is asserted without proof and attributed to (4.4), which is not
its source. The gap is load-bearing --- with the naive bound $6\alpha x+6$ the dimensions $L_a$
go negative at $x=3$ and Proposition 4.1 collapses --- but it is small: we supply the missing
two lines in Finding~\ref{f:ba}, and with them the argument goes through. By the ordinary
standard, a proof with an unproved load-bearing step is not correct, so the paper needs that
correction before it can be called one. It is a repairable defect, not a false statement and
not an unproved assumption carrying the theorem.

\textbf{With that gap filled, we found no further error.} We read every proof in the paper and
reconstructed the ones that carry weight; we found each of them complete. We could not, however,
\emph{certify} \S4.1: two careful readings of a dense $p$-adic valuation argument are evidence,
not refereeing, and unlike every other estimate in the paper its conclusion cannot be
cross-checked against a computed determinant, because its own hypothesis (4.1) forces a matrix
of size $\ge296000$. That is a limitation of our cross-checking, not a defect in the paper's
argument, and we say so plainly here because an earlier draft of this report phrased it in a way
that suggested otherwise.

So the verdict is: \emph{one repairable gap, which we have repaired; no other error found; and
one section whose proof we have read and believe but cannot certify.} What this is not is a
certification. ``No error found'' is not ``no error'', and the distinction matters most where we
could not compute --- \emph{if} \S4 over-claims $\gin$, the \emph{proof} fails while the
\emph{theorem} may well survive, for the reason given in Section~\ref{sec:musthold}.
\end{verdictbox}

\subsection*{Confidence, and what would change it}

We separate three questions, because our confidence in them differs.

\begin{enumerate}
\item \emph{Do Sections 1--3, 5--7 and Appendices A--C contain an error?}
\textbf{Very high confidence that they do not} ($\approx0.95$). Everything in
this range was verified independently, most of it twice by different methods; Appendix~A's
certification of Lemma 6.1, which the audit initially recorded as uncheckable, was in the end
fully reproduced (Section~\ref{sec:appA}), and Appendix~B reproduces exactly as rationals.
The figure is set by the weakest component rather than the strongest: the verifying pass on
\S5 and Appendix~B put itself at $0.97$ (conditional on (3.12), (4.8) and (4.14)), the pass on
\S6 and Appendix~A at $0.95$, and a conjunction cannot beat its weakest term. What is left in
this range is reading, not computing --- the proofs of Lemmas 3.1--3.3, 4.2 and 6.2 and the
partial summation in Proposition 5.2 were audited line by line and machine-tested only through
their consequences.

\item \emph{Is Section 4 correct?} \textbf{High confidence} ($\approx0.9$). The outer range
(Proposition 4.3) is verified against exact determinants at $37$ primes with \emph{equality at
27 of them}, which is very strong evidence for a formula of this complexity. The inner range
(Proposition 4.1) is proved in the paper (pp.~10--11), and we reconstructed that proof line by
line, twice and independently, and found it complete once the missing inequality of
Finding~\ref{f:ba} is inserted. What we could not do is check its conclusion against a computed
determinant, since (4.1) forces $p>8000$ and a matrix of size $\ge296000$; instead its mechanism
was tested at the smallest nondegenerate instance ($K=80$, $p=23$), where it reproduces the true valuation $v_{23}^{\mathrm{G}}(\Delta_{80})=-188$
\emph{exactly}, and where switching off the paper's side conditions makes the bound false. All
side conditions were swept over tens of thousands of primes at the real parameters with no
failure.

\item \emph{Is the theorem true?} \textbf{High confidence} ($\approx0.95$), and this is
worth separating from (2), because the evidence is of a different kind. By the paper's own
remark C.3, irrationality follows from
$\limsup K^{-2}\log P_K(\z)<0$ for the \emph{primitive} polynomials
$P_K=\Delta_K/\operatorname{cont}(\Delta_K)$, with no reference to $m_{K,M}$ at all. Measured
directly from exact determinants, $K^{-2}\log P_K(\z)=-0.1657,-0.1302,-0.1253$ at $K=40,80,120$,
against a requirement of $<-0.0174$: the construction has roughly seven times the room it needs.
An error in Sections 3--5 would therefore damage the \emph{proof} without necessarily touching
the \emph{theorem}.
\end{enumerate}

\noindent Overall we would put the probability that the paper is a correct proof of Theorem 1.1
at roughly $0.85$, and we would not be surprised to be told that a referee with more time found
a repairable gap in Section~4. We would be surprised to be told the paper is wrong.

\medskip\noindent\textbf{What would change the verdict.} (i) A demonstration that
$\vG(\Delta_K)<\gin$ for some admissible inner prime --- this is the one inequality nobody has
tested at the paper's parameters, and it would be fatal. (ii) An error in Lemma 3.1 or 3.2 of
one power of $p$, which would move $\gin$ by about $h$ and change $A_M$. (iii) Any correction to
$A_{200}$ or $\Ub$ exceeding $1.74\times10^{-2}$ ($1.29\%$ of $A_{200}$), which is what it takes
to reach Theorem 1.1. Three smaller thresholds sit below it and touch only the printed
constants: $5.07\times10^{-6}$ breaks (2.8) and with it Corollary 1.2, $3.28\times10^{-5}$ forces
the printed exponent $139/5$ in (2.7) to be reprinted, and neither disturbs Theorem 1.1.
Nothing we found moves either constant at all.

\medskip\noindent\textbf{On the prior.} The problem is old and famous and the author is
unknown to us; the paper carries a disclosure that a generative AI tool was used for
``consistency checks of the arguments, calculations, and references''. None of this is an
argument. We record that the paper passes a test that no wrong paper in this circle of ideas
has passed in our experience: its machinery reproduces $\zeta(3)$ (with a weak exponent
$\approx14.5$) and \emph{fails} for $\zeta(7)$, $\zeta(9)$, $\zeta(11)$, with balances
$+0.93$, $+1.97$, $+3.04$ against $-0.049$ for $\zeta(5)$. It does not prove too much.

\clearpage

\section{Summary table of the numbered statements}

Verdicts: \OK\ = checked and correct; \IM\ = correct but imprecise, terse or
not reproducible as written; \GP\ = a gap in the written argument; \FA\ = false; \NC\ = could
not check. ``Exact'' means exact rational or $p$-adic arithmetic.

\begingroup
\small
\setlength{\tabcolsep}{4pt}
\begin{longtable}{@{}p{2.55cm}p{1.75cm}p{9.5cm}@{}}
\toprule
\textbf{Statement} & \textbf{Verdict} & \textbf{Comment} \\
\midrule
\endfirsthead
\toprule
\textbf{Statement} & \textbf{Verdict} & \textbf{Comment} \\
\midrule
\endhead
\midrule
\multicolumn{3}{r@{}}{\emph{continued on the next page}}\\
\endfoot
\bottomrule
\endlastfoot

\multicolumn{3}{@{}l}{\textbf{Section 1}}\\
Theorem 1.1 & \OK & Deduction from Theorem 2.1 is airtight: quadratic decay beats linear height growth, so no height estimate is needed. Its truth is that of Theorem 2.1. \\
Corollary 1.2 & \IM & Argument valid, constants exact. But $b_0>\exp(1.07\times10^{11})$: the \emph{threshold} is not merely unspecified but ineffective, so the exponent 260 is not usable for any $b$ one could write down. The paper says only ``not specified''. \\
$0<|\z-p/q|<q^{-0.86}$ & \OK & Brown--Zudilin citation verified against arXiv:2210.03391 (lines 86, 180). \\
\addlinespace
\multicolumn{3}{@{}l}{\textbf{Section 2}}\\
(2.1) parameters & \OK & $1+2\alpha=23/20$, $h=K-N=37n$, $\lambda=37/40$, $\alpha=3/40$ all consistent. \\
(2.2) $\mu_X(t^e)$ & \OK & Symbolically identical to $(2e{+}5)!\,\zeta(2e{+}2)/(12(2\pi)^{2e+2})$ for $e=0..11$; matches $\int y^{2e}w$ to 24--40 digits. \\
(2.3) $\mu_X(1/(t{+}j^2))$ & \OK & Equals $a^4\zeta(5,a)-\tfrac1{2a}-\tfrac14$ at $a=j$; both correction terms and both signs verified, symbolically and to $\sim$24 digits, $j=1..6$. \\
\S2.1 extension of $\mu_X$ & \IM & No circularity: $\{t^e\}\cup\{1/(t+j^2)\}$ is a basis, so (2.2)--(2.3) \emph{define} rather than extend. Exclusion of $t=0$ is used but never stated. \\
(2.4)--(2.6) & \OK & $G_K$ genuinely Hankel; $S_K>0$; $m_{K,M}>0$ by construction. $S_K$ is independently forced by (3.11). \\
Theorem 2.1 & \NC & Degree and positivity verified. Integrality is Proposition 5.1; decay is (7.2). \\
(2.7), (2.8) & \OK & Follow from (7.2). Both constants are saturated: see Section~\ref{sec:numerics}. \\
(2.9) degree & \OK & Exact identity, sign included, at $(N,K)=(1,3),(2,5),(3,7),(1,5),(3,9)$ and at the real $(3,40)$, where it matches a 1594-digit integer. Hence $\deg Q_{K,M}=h$ exactly. \\
Proposition 2.2 & \OK & Every step checked: $w=y^5f^{(4)}/12$, $w(0^+)=1/\pi$ to 20 digits; the quartic derivative identity symbolically; all four boundary terms $O(y)$; Hermite (DLMF 25.11.29) to $\ge24$ digits at five values of $a$. Only omission: the $\sum/\int$ interchange (Tonelli, one line). \\
(2.10) & \OK & Positive definiteness of $G_K(\z)$ confirmed at $K=40,80,120,160,200,240$; at $K=40$, all 37 LDL pivots positive, certified five ways including interval arithmetic. \\
\addlinespace
\multicolumn{3}{@{}l}{\textbf{Section 3}}\\
\S3.1 defs, $\kappa_d$, $\|\tau\|$ & \OK & $\|\tau\|\le p$ is the correct power and is sharp (attained at $z^{4p-1}$). $B_1=-1/2$ \emph{is} stated (p.~3). \\
(3.1) pullback & \OK & Exact on monomials ($e=0..10$) and on poles ($j=1..12$); reproduces (2.2) and (2.3). \\
(3.2), (3.3) & \OK & From $B_n(1)=(-1)^nB_n$ and $B_n(x{+}1)-B_n(x)=nx^{n-1}$. The hypothesis ``$g$ regular at zero'' is necessary and consistent. \\
Lemma 3.1, (3.4) & \OK & Valid; tested exactly at $p=7,11,13$. $d(r_\nu)<p$ is necessary ($r=p$ gives $v_p=-5$); $\deg U_0\le p+1$ is sufficient but not tight (true threshold $4p-2$); $r_\nu-r_\eta\in\mathbb{Z}_p^\times$ is redundant given $d(r_\nu)<p$ (now proved, not just searched). \\
(3.5) & \OK & Exact; converges in the Tate algebra iff $v_p(s)<0$. \\
far/near consistency (p.~7) & \OK & True, and the paper's one-sentence justification (uniqueness of the $B_T$ decomposition) is the right one. Verified on 18 random instances. \\
Lemma 3.2, (3.6), (3.7) & \OK & $p^{-4}$ is right ($p^{-1}$ Raabe, $p^{-3}$ three derivatives); $Y=p^5X+C_p$ is exactly what makes the $1/x$ case return $-X$. $C_p\in p^5\mathbb{Z}_p$ for every $p\ge7$ (a theorem, not a table); ``$p\ge7$'' is genuinely needed. \\
Lemma 3.3, (3.8)--(3.10) & \OK & Every estimate checks; (3.8) and (3.9) are sharp (min slack 0). (3.10) never violated in $180+216$ exact $(K,d,p)$ cases on two different grids. Residue sign left as ``$\pm$''; it is $(-1)^{K-r}$. \\
(3.11) & \OK & \emph{The key ``no lost power of $p$'' test.} An exact identity: $\det[(K!)^2\mu_X(f_if_j/D_K)]=S_K\Delta_K$, with $S_K$ precisely (2.5). Verified exactly at five $(K,N,h)$ and then proved. \\
(3.12) & \OK & Follows from (3.10) plus the ultrametric determinant bound; holds at every prime $<5K+200$ at $K=40,80,120$ with large slack. \\
\addlinespace
\multicolumn{3}{@{}l}{\textbf{Section 4}}\\
(4.1) & \OK & Implies $p>200M\ge8000$ and $p^2>5K$; the latter \emph{is} used, on p.~14 in \S5.1. \\
(4.2), (4.3) & \OK & Exact $p$-power extractions, not bounds. The ``$-4$'' and ``$+1$'' silently absorb the $p^{-4}$ prefactor of (3.7). \\
(4.4), ordering rule & \OK & $L_0+\sum_aL_a=h$ is an identity. The rule ``extras to the largest $\ell_K$'' is not cosmetic: it excludes the one bad case of the weight comparison. \\
$b_a\le6\alpha x+3$ (p.~10) & \GP & \textbf{True and sharp, but asserted without proof and misattributed to (4.4).} Needs a case split on $v_N=N\bmod p$. With the naive $6\alpha x+6$ the chain gives $T-b_a>-3/2$ at $x=3$ and $L_a\ge0$ fails. See Finding~\ref{f:ba}. \\
(4.5) basis & \OK & $\mathbb{Z}_p$-unimodular by CRT; verified numerically ($v_p(\det U)=0$). \\
(4.6)--(4.8) & \OK & Re-derived; $\gin\in\mathbb{Z}$; equals the true valuation exactly in the one computable nondegenerate instance. \\
Proposition 4.1 & \NC & Proof reconstructed line by line and found complete. All three side conditions are displayed on pp.~10--11 (not hidden). \emph{Not testable at its own parameters} ($h\ge296000$); the $m_N\ge1$ sub-range has no determinant test. \\
(4.9) & \OK & Reads $5N\le2p-2$ (not $5K$). All five conditions hold throughout $K/3<p\le K$; the paper asserts this on p.~11. \\
(4.10) rank bound & \OK & $\operatorname{rank}L=r_p$ \emph{exactly} at all 37 tested primes --- the bound is sharp. \\
(4.11), (4.12) & \OK & Weights re-derived; $\vG(A_{uv})\ge w_u+w_v$ with 0 failures in $\approx3\times10^5$ entries. \\
divided difference (p.~12) & \IM & Conclusion true, but needs the $-\tfrac14+\tfrac1{2j}$ tail of (2.3): over 2087 node pairs the congruence holds 2087/2087 with the tail and 0/2087 without. One displayed line missing. \\
Lemma 4.2, (4.13) & \OK & Proved independently; 6000 random trials, 0 violations, bound attained. The nonpositive \emph{half}-integer hypothesis is demonstrably necessary (fails for quarter-integers in 102/3000 trials). \\
(4.14), table p.~13 & \OK & Rebuilt from (4.11)--(4.12)+Lemma 4.2: 0 mismatches over 3125+14100 $(K,p)$ pairs, and every row of the cost table re-derived. \\
Proposition 4.3 & \OK & Verified against exact $\Delta_K$ at $K=40,80,120$: \emph{0 violations at all 37 primes, equality at 27}. Second assertion ($p>K$) exact at every prime up to $4h$. \\
\addlinespace
\multicolumn{3}{@{}l}{\textbf{Section 5}}\\
(5.1) & \OK & The three branches partition the primes with no gap or overlap; boundary cases $pM=K$, $3p=K$, $p=2h$, $p>2h$ all check. Branch 1 bounds $F_K$, branches 2--3 bound $\Delta_K$ --- load-bearing, and correct. \\
(5.2), (5.3) & \OK & Only $p\le2h$ enters; (5.3) is exactly Legendre applied to (2.5), signs included. \\
Proposition 5.1 & \NC & Correct bookkeeping of (3.12), 4.1, 4.3. Inherits Proposition 4.1's status. \\
(5.4)--(5.6) $\Gamma,\mathcal{N},R$ & \OK & Re-derived from (4.4)--(4.8) and (5.3) independently; $R$ affine on exactly 143 intervals of $[3,20]$, matching the paper's $s=143$. \\
(5.7) uniformity & \IM & True and better than claimed: $|\gin-p\Gamma(K/p)|\le3.8M$, both signs, uniform in $K$ (exactly 35 at fixed $x$ over a factor 32 in $K$). The paper's one-paragraph justification is very terse and gives no constant; ``$O(M)$ intervals'' over-counts (there are at most 3). \\
(5.8), (5.9) & \OK & Re-derived from (4.14) as an exact $K\to\infty$ limit; 0 mismatches in 20000+30000 exact random-rational tests. $\int_{1/3}^{1/2}d=9/640$ exactly. \\
(5.10) $\Iout=127751/96000$ & \OK & Exact, by three independent routes; all 11 rows of Table 4 and their endpoints reproduce. \\
Proposition 5.2, (5.11) & \OK & PNT step correct: $p\le\sqrt{5K}$ gives $O(K^{3/2}\log K)$; $\lfloor\log_p5K\rfloor=1$ exactly on $(\sqrt{5K},K/M]$; partial summation legitimate. End-to-end control converges to the predicted blocks. \\
(5.12)--(5.14) & \OK & Re-derived by hand; the quadratic coefficient of $R$ cancels \emph{identically} in $(\alpha,\lambda)$ once $H=\lambda+3\alpha$, so convergence of $\int R/x^3$ is structural. $-1/2\le Q\le13/8$ correct. \\
(5.15)--(5.17) & \OK & Both integrations by parts re-derived; $\overline{\mathcal P}=2923/240$, $\max|G_0|=1/(36\sqrt3)$, $|\mathcal{C}|\le118511/(4320\sqrt3)<16$. (5.16) is an exact identity of its constants ($=-2689/48000$) and is conservative by $4.05\times10^{-3}$; (5.17) holds at $M=200$ with slack $2.0\times10^{-5}$. \\
(5.18)--(5.21) & \OK & $\int_3^{20}R/x^3$ exact; $A_*$, $A_{200}$, $A_{100000}$ reproduce exactly. \\
\addlinespace
\multicolumn{3}{@{}l}{\textbf{Section 6}}\\
(6.1), (6.3) $V$, $C_*$ & \OK & $V$ is exactly the field produced by $y=K\sqrt t$; $C_*=2.65303599034048866$ is exactly the $K^2$ coefficient of $\log S_K$ (symbolic difference 0). \\
Lemma 6.1, (6.2) & \OK & \emph{The crux, and it is both true and proved.} $\sup_{t\ge0}(2U^\rho-V)=-6.6498938806$ at $t^*=0.6195486669$, against $M_0=-6.645$: slack $4.894\times10^{-3}$. Independently reproduced here. \\
Lemma 6.1, (6.4) & \OK & $\lambda M_0-I(\rho)+C_*=-1.366995564512461\le\Ub=-1.3669955$. True, but the margin is $6.45\times10^{-8}$: $\Ub$ is literally the 7-decimal rounding of the left side. \\
Lemma 6.1, structure & \OK & $\sum c_j=37/40$ \emph{exactly} in integer arithmetic; nesting holds; $\min_j(b_j-a_j)=0.005085947445>1/225$. \\
Lemma 6.2, (6.5) & \OK & Standard and correctly proved; the integrability hypothesis is exactly what the circle/arcsine application needs. \\
(6.6)--(6.9) & \OK & Re-derived; arcsine window bound $4\sqrt{2r}/(\pi\sqrt L)$ correct; $\sum_j 8c_j\sqrt2/(\pi\sqrt{b_j-a_j})=9.42\le60$. Tested end-to-end on 180 random configurations at $K=40,80,200$. \\
(6.8) & \OK & True and sharp: the coefficient $13/10=2\lambda-1+6\alpha$ is exact and $6\alpha^3/t$ is exactly what $6\alpha\log(1+\alpha^2/t)\le6\alpha\cdot\alpha^2/t$ gives (an earlier pass called it $3\times$ too generous; that was a misreading). \\
(6.10) Andr\'eief & \OK & Verified symbolically and exactly for $h=2,3,4$, including the $1/h!$; weight identity checked numerically. \\
(6.11), (6.12) & \OK & $8192$ valid and generous (the sharp constant is $(2\pi)^4/12=129.88$). (6.12) correct, with the $t=0$ case the binding one. \\
(6.13)--(6.15) & \OK & Exponent $2h(h{+}6N{-}K)+16h$ and prefactor $4096e^{12}$ re-derived; $\int_0^\infty t^{-1/2}(1+\sqrt t)^5e^{-\sqrt t}dt=652$ exactly; the $K^2\log K$ terms cancel exactly ($111/160-111/160=0$). \\
Proposition 6.3, (6.16) & \OK & Holds at $K=40,\dots,240$ (legitimate instances, since the claim is for every $K\in40\mathbb{Z}_{>0}$) with room to spare (the excess $\log F_K-\Ub K^2$ peaks at $+276$ at $K=120$ and is $+27$ at $K=240$, against an allowance $24K\log K+200K$ rising from $11541$ to $79569$: a factor $79$ at $K=40$, $387$ at $K=200$, $2988$ at $K=240$). Note its own bound is vacuous for $K\le242$. \\
\addlinespace
\multicolumn{3}{@{}l}{\textbf{Section 7}}\\
(7.1) & \OK & A legitimate sum of two upper bounds; $Q_{K,M}(\z)>0$ so the logarithm exists. \\
(7.2) & \OK & Both exact margins reproduce digit for digit from Appendix B.3. \\
\addlinespace
\multicolumn{3}{@{}l}{\textbf{Appendix A}}\\
Table 1 & \OK & Transcribed from the text layer and checksummed: $\sum c_j=37/40$ exactly. \\
(A.1), (A.2) & \OK & (A.1) against quadrature to $<10^{-40}$; (A.2) re-derived from nesting and confirmed by double quadrature. \\
(A.3), (A.4) & \OK & Truncation orders ample: remainders $4.9\times10^{-64}$ and $2.1\times10^{-51}$, far below the claimed $2^{-144}$ and $10^{37}$ times below the tightest margin in play. \\
(A.5), min of $V$ & \OK & Matches direct integration of (6.1); $\Phi''$ sign change at $y^2=711/880$ exactly; minimum inside $(q_-,q_+)$ as claimed. \\
(A.6) & \IM & $V_*=0.913559526\le\min V=0.913559534$: valid. ``The parenthesis is negative'' holds by $2.5\times10^{-8}$ but follows in one line from the preceding sentence; terse, not unjustified. \\
(A.7), (A.8), Table 2 & \OK & The $k$-column lists whitespace-separated runs. Under that reading the 91 rows give \textbf{684 intervals that tile $[0,2]$ exactly} (verified here in exact rational arithmetic). The caption should say so. \\
(A.9) & \OK & $\max\mathcal{B}(l,r)=-6.645002689215588$ over all 684 intervals, below $-6645002/10^6$. \textbf{Margin $6.892\times10^{-7}$}; exactly one interval is within $10^{-5}$ of the threshold. \\
(A.10), (A.11) & \OK & Both enclosures correct; $\kappa=13.8658566788652898$ inside (A.11) and $<6933/500$. \\
\addlinespace
\multicolumn{3}{@{}l}{\textbf{Appendix B}}\\
(B.1) & \OK & Re-derived from $\ell(x,z)=\lfloor x-z\rfloor+\lfloor x+z\rfloor+1$. \\
(B.2) breakpoints & \IM & Correct and complete (142 interior points, $s=143$), but redundant: $2(1-\alpha)=2\lambda$ is listed twice and the $2\alpha$-points are a subset of the $4\alpha$-points. \\
(B.3), Tables 3, 4 & \OK & All 17 rows of Table 3 and all 11 rows of Table 4 reproduce exactly; the 17 rows sum \emph{exactly} to (5.18). (See Section~\ref{sec:disagree} on a disputed digit.) \\
B.3 final margins & \OK & Both printed rationals reproduce exactly as fractions. \\
\addlinespace
\multicolumn{3}{@{}l}{\textbf{Appendix C}}\\
Proposition C.1, (C.1) & \OK & (C.2)--(C.5) and the combination re-derived; $\kappa<6933/500$ with margin $1.433\times10^{-4}$. The full chain was run numerically at $K=40$ and $80$: (C.3) holds as an exact identity to 401 and 1500 digits. \\
(C.4) & \OK & Uses $m!\ge(m/e)^m$ and monotonicity of $x\log(eK/(\alpha x))$; valid since $h=37K/40<40K/3$. \\
(C.5) & \IM & The displayed $h^2\mathcal{R}^{2h}$ \emph{is} reproducible from the paper's own argument and is true; it is weaker by the $O(1)$ factor $\mathcal{R}^2=293.8$ than the $h^2\mathcal{R}^{2(h-1)}$ the same step gives. $\kappa$ unaffected. \\
(C.6) & \OK & $\kappa=13.86585667886528983$ by two evaluators agreeing to $>100$ digits; margin to $\varrho$ is $1.4332\times10^{-4}$. The ``positive margin absorbs the polynomial factor'' only for $K\gtrsim2.81\times10^5$; neither number is given. \\
(C.7), C.2 & \OK & All rational identities exact: $\varrho c-260=-1/80$, $\epsilon-\lambda/c=1/24000$. The perturbation argument genuinely excludes a zero. Depends on Proposition 5.1 as well. \\
(C.8) & \OK & $\lambda\varrho-\epsilon=511067/40000=12.776675<13$ exactly. \\
C.3 primitive normalization & \OK & Gauss's lemma; and it is the remark that makes the theorem robust to an error in \S\S3--5 (see Section~\ref{sec:musthold}). \\
\end{longtable}
\endgroup

\section{Per-section commentary}

\subsection{Sections 1 and 2: the construction}

Nothing here is wrong, and the parts a referee would most want to see are the parts that are
most solid.

The final logical step is airtight and deserves to be stated plainly, because it is what makes
the paper's strategy different from a linear-form construction: one needs only
$Q_n\in\mathbb{Z}[X]$, $\deg Q_n\le37n$, $Q_n(\z)>0$ and $Q_n(\z)<e^{-139n^2/5}$. Then
$b^{37n}Q_n(a/b)$ is a positive integer smaller than $\exp(37n\log b-\tfrac{139}{5}n^2)\to0$.
The decay is quadratic in $n$ and the height cost is linear, so \emph{no estimate on the
coefficients of $Q_n$ is required at all}. This is why the paper can afford to work with
polynomials of growing degree.

The functional of (2.2)--(2.3) is well defined and there is no circularity, contrary to what a
quick reading might suggest. The set $\{t^e\}_{e\ge0}\cup\{1/(t+j^2)\}_{j\ge1}$ is a basis of
the space of rational functions with simple poles in $\{-j^2\}$, so (2.2) and (2.3) prescribe
values on a basis and determine a unique $\mathbb{Q}$-linear functional; the word ``extend'' in
the paper is loose for ``define''. One hypothesis is used and never stated: $0=-0^2$ is not an
allowed pole (all products run over $j\ge1$), which is exactly why the integral in
Proposition 2.2 converges.

Equation (2.9) is correct including the sign, and this matters: it is what gives
$\deg Q_{K,M}=h$ \emph{exactly}, which the final step needs. We verified it symbolically on
small cases in both sign parities, and then on the paper's own smallest real parameters
$(N,K)=(3,40)$, where $\det[X]G_{40}$ equals
$(-1)^{666}\prod_{j=4}^{40}j^4D_3(-j^2)^5$ exactly, as a 1594-digit integer, reproduced by two
tracks with different determinant routines. A structural fact the paper never states is worth
knowing: $[X]\mu_X(D_N^6t^m/D_K)$ is an $(h{-}1)$st divided difference and vanishes identically
for $m<22n-3$, so $[X]G_K$ carries a large zero block and the Vandermonde nonvanishing in (2.9)
is genuinely load-bearing rather than decorative.

Proposition 2.2 was checked step by step and is right, including both correction terms and both
signs in (2.3) --- the place where such a computation usually goes wrong. The four integrations
by parts were certified by comparing two different integrands to 45--51 digits; Hermite's
formula (DLMF 25.11.29) was matched to $\ge24$ digits at five values of $a$. Positive
definiteness of $G_K(\z)$ was confirmed at $K=40,80,120,160,200,240$, and at $K=40$ certified
five ways including interval arithmetic. A warning for anyone reproducing this: the entries
carry up to 107 digits of cancellation, and evaluating $\Delta_{40}$ at $\z$ needs $\ge2600$
digits --- at 400 digits it returns $+1203.8$, pure noise. The single omission in the proof is
the interchange of $\sum_{\ell\ge1}$ with $\int_0^\infty$, which is Tonelli in one line since
every term is positive.

\subsection{Section 3: the local functional}

\textbf{No error, and --- the thing we were most concerned to check --- no place where a power
of $p$ is gained or lost silently.} Some 400 automated checks in exact arithmetic with explicit
valuations, plus a SageMath/$\mathbb{Q}_p$ cross-check of the Bernoulli inputs, produced no
failure.

The decisive test is (3.11). A change of basis is exactly where a construction of this kind
loses a power of $p$ without anyone noticing, and here the change of basis to the
integer-valued $f_i$ is an \emph{exact identity}: with $q_i=\sum_kC_{ik}t^k$ triangular with
$C_{ii}=(-1)^i2/(2i)!$, one gets
$\det[(K!)^2\mu_X(f_if_j/D_K)]=((K!)^2/(N!)^{12})^h(\det C)^2\Delta_K=S_K\Delta_K=F_K$,
with $S_K$ \emph{precisely} the constant printed in (2.5). We verified this exactly for five
$(K,N,h)$ triples and then proved it. Nothing is gained or lost by the change of basis, and
nothing by the pullback $t=-x^2$.

Lemma 3.1's hypotheses were attacked individually. $d(r_\nu)<p$ is \emph{necessary}: at $r=p$
the conclusion fails with $v_p=-5$. $\deg U_0\le p+1$ is sufficient but far from tight --- the
true threshold is $4p-2$, because $\kappa_d=d(d-1)(d-2)B_{d-3}/24$ has a compensating factor of
$p$ at $d=p+2,2p+1,3p$, and the first loss is at $d=4p-1$ (measured: $27,43,51,67,75,91$ for
$p=7,11,13,17,19,23$). The third hypothesis, $r_\nu-r_\eta\in\mathbb{Z}_p^\times$, is
\emph{redundant} given $d(r_\nu)<p$: inside $\{-p,\dots,p-1\}$ a congruent pair must be
$\{r,r-p\}$, and $H^{(5)}_{p-1-r}\equiv H^{(5)}_r\pmod p$ makes the two $p^{-1}$-sized residue
contributions cancel. This is unused slack, not an error.

Lemma 3.2 is correct and the exponent $p^{-4}$ is right ($p^{-1}$ from Raabe's multiplication
theorem, $p^{-3}$ from three differentiations). $Y=p^5X+C_p$ is exactly the choice making the
right-hand side at $g=1/x$ equal $-X$. The hypothesis $p\ge7$ is genuinely required, not
decorative: $v_p(\sum_{a<p}a^{-4})=1$ for $p\ge7$ but $0$ for $p=5$. One refinement for the
record: $C_p\in p^5\mathbb{Z}_p$ for every $p\ge7$ is a theorem, but $v_p(C_p)=5$ \emph{exactly}
only when $p\nmid B_{p-5}$; at the irregular prime $p=37$ one has $v_{37}(C_{37})=6$. The paper
claims only membership, so it is unaffected.

The one defect in \S3 is cosmetic: in Lemma 3.3's proof the residue coefficient is written
$c_r=\pm(K!)^2rA(r)/((K+r)!(K-r)!)$ with the sign left to the reader. It is uniformly
$(-1)^{K-r}$, and only $|c_r|_p$ is used, so nothing depends on it.

\subsection{Section 4: the local determinant estimates}

This is where the risk lives, and it is worth being precise about \emph{which} risk.

\textbf{The outer range is verified, and it is essentially sharp.} We computed $\Delta_K(X)$
exactly in $\mathbb{Q}[X]$ at $K=40,80,120$ (coefficients up to $2.7\times10^4$ digits) by two
independent implementations, and compared $\vG(\Delta_K)$ with $\gout$ at every prime satisfying
(4.9). There are \textbf{0 violations at all 37 primes, with equality at 27 of them}. A formula
of this complexity that is exactly attained at 27 independent primes is almost certainly right.
The second assertion of Proposition 4.3 ($\vG(\Delta_K)\ge0$ for $p>K$) is exact at every prime
up to $4h$, which is nontrivial because the residue denominators contain primes up to $2K-1$
and they cancel. Every ingredient was separately verified: unimodularity of the basis (4.11)
at 37/37 primes; $\operatorname{rank}L=r_p$ \emph{exactly} at 37/37, so the rank bound (4.10) is
sharp; $\vG(A_{uv})\ge w_u+w_v$ with 0 failures in $\approx3\times10^5$ entries; and the
combinatorial identity (4.14) rebuilt from (4.12) and Lemma 4.2 with 0 mismatches over more than
$1.7\times10^4$ $(K,p)$ pairs. Where $\gout$ is not sharp --- at $K\ge2p$, where it misses by
$78,60,17$ at $K=120$ --- the loss is in the \emph{safe} direction: it inflates $A_M$ and makes
(7.2) harder. Section~4 therefore cannot rescue a false theorem, which is a useful thing to
know: if the paper is wrong, it is not wrong here.

Lemma 4.2 we proved independently and stress-tested (6000 random instances over
$\mathbb{Z}_7$, $\mathbb{Z}_{11}$, $\mathbb{Z}_3$; 0 violations; the bound attained in every
sharpness test). Its restriction to nonpositive \emph{half}-integer weights is not a
convenience: with quarter-integer weights the conclusion fails in 102 of 3000 random
instances, e.g.\ $p=11$, $h=3$, $w=(-3/2,-1/4,-3/2)$, $r=1$, $z=0$, where the claimed bound is
$-13/2$ and the true valuation is $-7$. The paper says the restriction is essential and it is
right.

\textbf{The inner range: the proof was read and reconstructed; only the numerical cross-check
is missing.} Proposition 4.1 is proved in the paper on pp.~10--11 --- about twenty-five lines,
resting on Lemmas 3.1 and 3.2 and on the allocation (4.4)--(4.8), with the side conditions
displayed in the text. We reconstructed that argument line by line, twice and independently, and
found it complete apart from the one asserted inequality of Finding~\ref{f:ba}, which we prove
below. We record two things about our confidence in it. First, it is a reading, not a
verification: it is the one substantial argument in the paper whose correctness rests on our
having followed it rather than on a machine check. Second, and unlike every other estimate here,
its conclusion cannot be compared against a computed determinant: hypothesis (4.1) forces
$p>200M\ge8000$ and $K\ge200M^2\ge320000$, hence $h\ge296000$, and at the computable $K\le160$ the
allocation (4.4) is degenerate ($L_0=4M+10>h$, and $T<0$). So the corroboration we can give it is
indirect:

\begin{itemize}
\item \emph{The mechanism, at the smallest nondegenerate instance.} At $K=80$, $p=23$ with $L_0$
freed, the true value is $v_{23}^{\mathrm{G}}(\Delta_{80})=-188$. With $L_0=3$ or $4$ (all side
conditions satisfied) the entry inequality holds with \textbf{0 failures} and $\gin=-188$
\emph{exactly}. With $L_0=0,1,2$ (zero-source dominance fails) the entry inequality fails
$2094/812/29$ times and $\gin$ would be $-170/-180/-186$, i.e.\ a \emph{false} bound. The test
is therefore discriminating, and the paper's side conditions are exactly the right ones.
\item \emph{The side conditions, at the real parameters.} All of them --- $L_a\ge0$, the two
Lemma-3.1 degree restrictions, zero-source dominance, the dimension identity
$L_0+\sum_aL_a=h$, and $2Hx-21/20<T<2Hx$ --- were swept over all 9158 inner primes at
$M=40$, $K=200M^2$, and over $27171+13960$ sampled primes at $M=200$, $K=8\times10^6$ and
$4\times10^7$. \textbf{No failures}, with three to four orders of magnitude of slack in the
degree conditions.
\item \emph{The hand-off from \S3.} The sharpest worry raised in the audit was whether \S4
actually produces an honest \emph{polynomial} $U_0$ of degree $\le p+1$ for Lemma 3.1, given
that the raw pullback numerator has degree $184n$. It does, and the paper displays exactly this
on p.~11: $\nu_i(c)+\nu_j(c)+6\ell_N(c)\le2(T+1)<\tfrac{23}{5}M+2$ at an ordinary source and
$5+4L_0+12m_N\le\tfrac{169}{10}M+45$ at the zero source, both $\le p+1$ since $p>200M$. We
verified both over the same sweeps. This closes the hand-off.
\end{itemize}

\noindent What remains without an independent machine check is the $p$-adic determinant
inequality $\vG(\Delta_K)\ge\gin$ itself at admissible parameters, and with it the $m_N\ge1$
(i.e.\ $p\le N$) part of the inner range, which for $M=200$ is most of it. To be clear about what
that does and does not mean: the paper proves this inequality, and we believe the proof; what we
cannot do is confirm it the way we confirmed Proposition 4.3, by computing the determinant and
comparing. Certification here needs a referee or a formalisation, and \S\ref{sec:musthold} says
so in those terms.

\begin{finding}\label{f:ba}
\textbf{(p.~10, the line after (4.4); the only genuine gap in the written mathematics.)}
The paper writes ``From (4.4), $2Hx-\tfrac{21}{20}<T<2Hx$, $b_a\le6\alpha x+3$.''
The attribution is wrong --- $b_a\le6\alpha x+3$ is a fact about $\ell_N$, not a consequence of
(4.4) --- and no proof is given.
\end{finding}

\noindent\emph{Why it matters.} The naive estimate $\ell_N(a)\le2\lfloor N/p\rfloor+2$ gives only
$b_a\le6\alpha x+6$, under which the paper's own chain reads
$T-b_a>2\lambda x-\tfrac{141}{20}$, which at $x=3$ equals $-3/2$: the nonnegativity of the
dimensions $L_a$ fails and Proposition 4.1 collapses. \emph{Why it is true.} Writing
$\ell_N(a)=2m_N+\mathbf{1}_{a\le v_N}+\mathbf{1}_{a\ge p-v_N}$ with $v_N=N\bmod p$, both
indicators can be $1$ only if $v_N\ge(p+1)/2$, and in that case $\alpha x-m_N=v_N/p>1/2$; so
$\ell_N(a)<2\alpha x+1$ either way. \emph{Evidence.} We checked the inequality over all 9158
inner primes at $M=40$, $K=320000$ (and 17214 at $K=640000$, 190197 at $M=200$,
$K=8\times10^6$): no violation, and the bound is \emph{strict everywhere} but only just ---
the slack is $6v_N/p$ when $1\le v_N\le(p-1)/2$ and $6v_N/p-3$ when $v_N\ge(p+1)/2$, so it
becomes arbitrarily small as $v_N/p$ approaches $0$ or $1/2$. The observed minimum is
$0.00119$, at $p=47981$, $v_N=24000$ (an earlier track recorded ``equality whenever $p\mid N$'';
that is wrong --- at $p\mid N$ the slack is exactly $3$). With the paper's bound the chain gives
$T-b_a>2\lambda\cdot3-\tfrac{81}{20}=+3/2$ exactly, so there is no room to spare at $x=3$ ---
this is not a place to be terse. The author should be asked to insert the case split.

\subsection{Section 5: normalization and the prime sum}

Section 5 and Appendix B are \emph{correct}, conditional on (3.12), (4.8) and (4.14). Every item
reproduced exactly in independent implementations. The case analysis of (5.1) is exhaustive and
non-overlapping, and one subtlety is worth recording because it is load-bearing and right:
branch 1 bounds $\vG(F_K)$ directly via (3.12) and therefore correctly \emph{omits} $v_p(S_K)$
from $L_p$, whereas branches 2--3 bound $\vG(\Delta_K)$ and must include it. Since
$v_p(S_K)\sim-1110n^2/(p-1)<0$ for small $p$, branch 1 would be false if (3.12) were a statement
about $\Delta_K$ rather than $F_K$. It is about $F_K$ (p.~9).

The single weak point of the exposition is (5.7). The claim that
$\gin=p\Gamma(K/p)+O_M(1)$ uniformly in the inner range carries the entire arithmetic constant
and is justified in one informal paragraph with no explicit constant and no displayed estimate.
The claim is \emph{true}, and in fact better than the paper suggests: computing $\gin$ exactly
from (4.4)--(4.8) over the full inner range gives
$|\gin-p\Gamma(K/p)|\le3.8M$, taking both signs, uniform in $K$ --- at fixed $x=K/p\approx10$
the difference is \emph{exactly} 35 for $K$ ranging over a factor of 32. So the error is $O(M)$,
not $O(M^2)$. The mechanism is a cancellation rather than an omission: the reserved zero block
(4.7) is genuinely $\Theta(M^2)$ (it contributes $+13734$ at $p=8081$, $K=320000$) but is almost
exactly cancelled by the ordinary block because $L_0=4M+10$ enters (4.4) and lowers $T$. That is
precisely what the paper's sentence about ``altering the number of allocated rows by $O_M(1)$''
means, but a reader cannot see it. A minor slip in the same paragraph: ``There are $O(M)$ such
intervals'' over-counts --- in the $z$ variable the integrand of (5.4) is a step function with
at most three pieces. The over-count is in the safe direction.

Everything else in \S5 is exact. $R$ is affine on exactly 143 intervals of $[3,20]$, matching
the paper's $s=143$, with zero affineness failures on any of them; $\int_3^{20}R/x^3$,
$\Iout=127751/96000$, $\int_{1/3}^{1/2}d=9/640$, $A_*$, $A_{200}$ and $A_{100000}$ all reproduce
as exact rationals by more than one route. The tail bounds (5.16) and (5.17) are valid and
conservative: the true $\int_{20}^\infty R/x^3=-0.06007376$ against the printed
$-2689/48000=-0.05602083$, so the paper discards $4.05\times10^{-3}$ there, which is 23\% of its
own final margin.

\subsection{Section 6: the real bound}

Correct, and --- this is the part of the audit that changed most --- \emph{proved}, not merely
true. Our initial reading recorded Lemma 6.1 as ``could not check'' because Appendix A's
certification looked unverifiable; it is verifiable, and it verifies. See
Section~\ref{sec:appA}.

Lemma 6.2 is a correct proof of a standard fact, with the integrability hypothesis exactly
matched to the circle-plus-arcsine application. It is applied correctly: $\sigma=K^{-1}\sum\omega_i$
and $\rho$ are both compactly supported and $\sigma-\rho$ has mass
$h/K-\lambda=37n/40n-37/40=0$ exactly. The regularization bookkeeping in (6.6)--(6.9) is right
in every direction, and the constant 60 is honest (the true value of
$\sum_j8c_j\sqrt2/(\pi\sqrt{b_j-a_j})$ is $9.42$; even the crude route using only
$b_j-a_j>1/225$ gives $49.97$). The hypothesis $b_j-a_j>1/225$ is used in exactly one place, here.

The scaling deserves a note for anyone reading the paper: it is $y_i=K\sqrt{t_i}$, i.e.\
$y^2=K^2t$, and only that scaling makes the bookkeeping in (6.13) close. We re-derived that
exponent tracking every power of $K$ --- Vandermonde $K^{2h(h-1)}$, and per variable
$K^{12N}\cdot K^{-2K}\cdot K^{12}\cdot K^5\cdot K^1$ --- and obtained
$2h(h+6N-K)+16h$, exactly as printed, with prefactor $8192/2\cdot e^{12}=4096e^{12}$, also as
printed. The integral $\int_0^\infty t^{-1/2}(1+\sqrt t)^5e^{-\sqrt t}dt=652$ is exact
($=2e\,\Gamma(6,1)$). The $K^2\log K$ terms of (6.14) and (6.15) cancel exactly, and the
$K^2$ coefficient of (6.15) equals $C_*$ of (6.3) with symbolic difference zero. Andr\'eief's
identity (6.10) is exact for $h=2,3,4$ including the $1/h!$.

\textbf{How much room \S6 actually has.} Four tracks independently solved the underlying weighted
equilibrium problem --- by a constrained convex QP with exact cell energies, by one-cut resolvent
inversion, by a Fourier--cosine Frostman solve, and by Frank--Wolfe on a log-spaced grid, each
with controls reproducing the semicircle, Marchenko--Pastur and quartic-field laws --- and they
agree to ten digits: the true optimum is
$E^*=\sup_\mu(I(\mu)-\int V\,d\mu)=-4.0385457735$, on $\mathrm{supp}=[7.12\times10^{-5},0.753950]$
with Robin constant $-6.6607811230$. The paper's comparison bound $\lambda M_0-I(\rho)=-4.0200316$
is \emph{above} $E^*$, hence valid, and loose by $0.01851$. In other words the true constant is
$\Ub_{\mathrm{true}}=E^*+C_*=-1.3855098$ against the paper's $\Ub=-1.3669955$: \textbf{Section 6
holds $0.0185$ of unused reserve, slightly more than the entire final margin
$|A_{200}+\Ub|=0.0174$.} The real side is conservative, not optimistic. This is confirmed from
the other direction by direct measurement: $K^{-2}\log F_K(\z)$ at $K=40,\dots,200$ extrapolates
to $-1.3854\ldots-1.3869$, and a model fitted on $K\le160$ predicts the $K=200$ value with an
out-of-sample error of $1.4\times10^{-6}$.

The consequence for the audit is sharp: \emph{any failure of this paper must be located in
Sections 3--5, not in Section 6.} If the arithmetic side were off by as much as $+0.036$, the
proof could in principle still be rescued by replacing $\rho$ with the true equilibrium measure.

\subsection{Section 7}

Two lines, both correct. (7.1) is a legitimate sum of two upper bounds, and $Q_{K,M}(\z)>0$ so
the logarithm exists. (7.2) reproduces digit for digit. The only observation is that (7.1) is a
$\limsup$ statement while (2.7)--(2.8) are ``for all sufficiently large $n$'': no threshold
anywhere in the paper is effective. That is legitimate for an irrationality proof, but it means
no finite computation can ever confirm or refute (2.7) directly --- every check must be of the
asymptotic constants, which is what this audit did.

\subsection{Appendix A}\label{sec:appA}

Appendix A proves Lemma 6.1, and it does prove it. Because the audit initially doubted this, we
state the reproduction in full.

Table 1 gives $10^{12}a_j$, $10^{12}b_j$, $10^{12}c_j$ as integers. Taken from the PDF text
layer (\emph{not} from the rendered page image --- see Section~\ref{sec:disagree}) it
checksums: $\sum c_j=37/40$ \emph{exactly} in integer arithmetic, the nesting
$0<a_{16}<\cdots<a_1<b_1<\cdots<b_{16}<2$ holds, and $\min_j(b_j-a_j)=0.005085947445>1/225$,
attained at $j=1$. From (A.2), $I(\rho)=-2.126593445147050403254$, inside the printed enclosure
(A.10); and $C_*=2.653035990340488661129$, also inside. Hence
\[
\lambda M_0-I(\rho)+C_*=-1.366995564512460935617\le\Ub=-\tfrac{2733991}{2000000},
\]
which is (6.4) --- \textbf{with a margin of $6.45\times10^{-8}$}. $\Ub$ is literally that number
rounded up at the seventh decimal.

The pointwise inequality (6.2) holds comfortably:
$\sup_{t\ge0}(2U^\rho(t)-V(t))=-6.6498938806$ at $t^*=0.6195486669$, against $M_0=-6.645$, a
slack of $4.894\times10^{-3}$; on $[2,10^8]$ the supremum is $-8.27$. We reproduced the
field $V$ of (A.5) against direct quadrature of (6.1) to $5\times10^{-51}$ as a control.

The certification itself, (A.7)--(A.9), is the part that looked unverifiable, because Table 2's
$k$-column contains rows like ``$19\ \ 3\ \ 0\ \ 2,3$'' and ``$21\ \ 5\ \ 2,\dots,4\ \ 10,\dots,13$''
whose format the caption does not explain. The rule is that the column lists a \emph{union of
whitespace-separated runs}. Under that reading all 91 rows parse, and the resulting
\textbf{684 intervals tile $[0,2]$ exactly} --- we checked this in exact rational arithmetic,
gap-free and overlap-free within each of the 35 blocks $[A_j,A_{j+1}]$. Evaluating $\mathcal{B}(l,r)$
of (A.7) on all 684 gives
\[
\max_{(j,d,k)}\mathcal{B}(l,r)=-6.645002689215588,\qquad
\text{at }(j,d,k)=(30,6,26),
\]
\[
[l,r]=[0.5480159423,\,0.5492641749],
\]
which is below the printed $-6645002/10^6$. \textbf{So (A.9) is true, with a margin of
$6.892\times10^{-7}$}, and exactly one of the 684 intervals lies within $10^{-5}$ of the
threshold. The prescribed truncation orders are ample: (A.3) at $m=64$ on $z\le1/3$ has
remainder $4.9\times10^{-64}$ and (A.4) at $m=80$ on $z\le1/2$ has remainder $2.1\times10^{-51}$,
both far below the claimed $2^{-144}\approx4.5\times10^{-44}$ and some $10^{37}$ below the
tightest margin in play. A margin of $6.9\times10^{-7}$ on quantities of size $6.6$ is resolved
by double precision, let alone by the 25--50 digits we used; the whole appendix verifies from
the printed tables in seconds (our own run: $4$\,s; the slowest of the four independent
reproductions took two minutes).

Two remarks. First, the author is operating at the edge of his own stated precision in
three places --- (6.4) at $6.45\times10^{-8}$, (A.9) at $6.89\times10^{-7}$ and (5.17) at
$2.0\times10^{-5}$ --- and one fewer dyadic refinement on the worst Table-2 cell would break
(A.9) as written. Second, and more seriously as a matter of robustness: the slack in (6.4) is
the same order as a single transposed digit pair inside one Table 1 entry. We confirmed this
experimentally: reading $10^{12}b_8$ as $201105672729$ instead of the correct $201105762729$
moves $I(\rho)$ by $3.3\times10^{-8}$. The printed table is self-consistent and correct, but a
referee is entitled to ask that (6.2) be quoted at $-6.6498$ (slack $4.9\times10^{-3}$) rather
than at $M_0=-6.645$. That single change replaces $\Ub$ by $-1.3715224$ and turns the margin in
(7.2) from $0.052484$ into $7.29553$ --- a factor of $139$ --- at no cost anywhere else. (We
recompute this here; the figure ``280'' that circulated in the audit's drafts was wrong.)

\subsection{Appendix B}

Exact throughout. All 17 rows of Table 3 and all 11 rows of Table 4 reproduce, the 11 row
endpoints coincide with independently derived breakpoints, and the 17 rows of Table 3 sum
\emph{exactly} to (5.18). (B.1) re-derives from $\ell(x,z)=\lfloor x-z\rfloor+\lfloor x+z\rfloor+1$.
The quadratic-cancellation identity
$2\lambda-12\lambda\alpha+2(H^2-H-\lambda^2)-18\alpha^2+6\alpha=0$
holds exactly, and in fact vanishes \emph{identically} in $(\alpha,\lambda)$ once one notes
$H=\lambda+3\alpha$ (which (4.4) forces but the paper never says): the convergence of
$\int R/x^3$ is structural, not a numerical accident of $3/40$ and $37/40$. Stating this would
make \S5.3 considerably more convincing. The breakpoint set (B.2) is correct and complete but
redundant: $2(1-\alpha)=2\lambda=37/20$ is listed twice, and the $2\alpha$-points are a proper
subset of the $4\alpha$-points.

\subsection{Appendix C}

Correct. Proposition C.1's chain (C.2)--(C.5) was re-derived and then run numerically end to end
at $K=40$ and $K=80$, where (C.3) holds as an exact identity to 401 and 1500 digits --- this
confirms the residue structure, the rescaling and the exponent $12N-2K+6$ simultaneously. The
constant is $\kappa=13.86585667886528983095$ by two evaluators agreeing to more than 100 digits,
against $\varrho=6933/500=13.866$: (C.6) is true with margin $1.4332\times10^{-4}$, and the
paper's own enclosure (A.11) is correct. All the rational identities are exact:
$\varrho c-260=-1/80$, $\epsilon-\lambda/c=1/24000$, $\lambda\varrho-\epsilon=511067/40000=12.776675<13$.
The perturbation argument in C.2 genuinely excludes a zero, and the exponent 260 is essentially
forced by this method (its best possible value is $259.54$).

Three caveats the paper should state. (i) ``Its positive margin absorbs the polynomial factor in
$h$'' needs $K\gtrsim2.81\times10^5$; neither the margin nor the threshold is given. (ii) The
true norm is much smaller than the bound: $\log\|B_K\|/K=7.314$ at $K=40$ and $7.388$ at $K=80$
against the asserted $\varrho=13.866$, so Proposition C.1 has roughly a factor-2 margin in the
exponent wherever it is testable, and the exponent 260 is an artefact of the estimate --- a
proved rate near $7.4$ would give about 139 by the identical argument. (iii) Corollary 1.2's
threshold is not merely unspecified but \emph{ineffective}: $M=100000$ forces
$K\ge200M^2=2\times10^{12}$, hence $\log b\ge1.067\times10^{11}$, i.e.\
$b_0>\exp(1.07\times10^{11})$. A larger $K$ is no escape, since $e^{\varrho K}$ then blows up.
Finally, Corollary 1.2 depends on Proposition 5.1 and on Theorem 2.1's degree claim as well as
on (7.1), Proposition 2.2, \S2.3 and (C.6) --- without integrality, $b^hQ_K(a/b)$ is not an
integer and the contradiction evaporates.

\section{The decisive numerics}\label{sec:numerics}

\subsection{The two constants and the balance}

Everything in Table~\ref{tab:balance} was recomputed here in exact rational arithmetic from the
paper's own formulas and matches the printed values.

\begin{table}[htbp]
\centering
\small
\begin{tabular}{@{}lll@{}}
\toprule
Quantity & Value & Status \\
\midrule
$\Iout$ & $127751/96000=1.3307395833$ & exact, three routes \\
$\int_3^{20}R\,x^{-3}dx$ & $0.0427881769084$ & exact, matches (5.18) \\
$A_*$ & $1.3175069269083709$ & exact, matches (5.19) \\
$A_{200}$ & $1.34958769774170422$ & exact, matches B.3 \\
$A_{100000}$ & $1.31757167571548622$ & exact, matches B.3 \\
$\Ub$ & $-2733991/2000000=-1.3669955$ & paper's constant \\
$\Ub_{\mathrm{true}}$ & $-1.3855098$ & true limit; $\Ub$ valid, loose by $0.0185$ \\
\midrule
$A_{200}+\Ub$ & $-0.017407802258296$ & $<0$: irrationality holds \\
relative margin & $1.2899\%$ of $A_{200}$ & \textbf{the margin the theorem turns on} \\
vs.\ $-139/8000$ & slack $3.2802\times10^{-5}$ & $0.0024\%$: the printed exponent is saturated \\
$A_{100000}+\Ub$ & $-0.049423824284514$ & $<-79/1600$, slack $4.882\times10^{-5}$ \\
$A_{200}+\Ub_{\mathrm{true}}$ & $-0.0359221$ & the margin with the sharp real constant \\
\bottomrule
\end{tabular}
\caption{The balance. Both B.3 margins reproduce exactly:
{\footnotesize$-1600(A_{200}+\Ub)-\tfrac{139}{5}=\tfrac{3089837638249482469}{58872425992515135000}$}
and {\footnotesize$-1600(A_{100000}+\Ub)-\tfrac{7907}{100}
=\tfrac{29873543950273155160680943}{3679526624532195937500000000}$}.}
\label{tab:balance}
\end{table}

Two sensitivities are worth separating, because they have different consequences.
\emph{Theorem 1.1} needs only $A_{200}+\Ub<0$ and has $1.29\%$ of room (and $1.60\%$ if one uses
(5.11) directly instead of the deliberately lossy (5.20)). The \emph{printed constants} $139/5$
and $7907/100$ survive perturbations of only $3.28\times10^{-5}$ and $5.07\times10^{-6}$; any
correction above those forces them to be reprinted and breaks Corollary 1.2, while leaving
Theorem 1.1 standing. Note also that $\Iout$ alone is $98.60\%$ of $A_{200}$: the arithmetic
constant is governed by (4.14) and Table 4 --- that is, by whether Proposition 4.3 and
Lemma 4.2 are true --- and not by any arithmetic Section 5 performs.

\subsection{The direct computations}

$\Delta_K(X)$ was built exactly in $\mathbb{Q}[X]$ from (2.1)--(2.6) alone, by four independent
implementations (multimodular determinant plus Lagrange interpolation; rational matrix inverse
plus characteristic polynomial; entry assembly from the closed residue formula plus Lagrange
interpolation over $\mathbb{Q}$; Berkowitz over $\mathbb{Q}[X]$), with controls passing first
(Hilbert determinants through the same code path; (2.2) against the moment integral; (2.3)
against $\int w/(y^2+j^2)$). All four agree coefficient by coefficient.

\begin{table}[htbp]
\centering
\begin{tabular}{@{}rrrrr@{}}
\toprule
$K$ & $\log F_K(\z)$ & $K^{-2}\log F_K(\z)$ & $\log(1/\mathrm{cont})$ & $K^{-2}\log P_K(\z)$ \\
\midrule
40  & $-2041.19514254$  & $-1.2757470$ & $1776.06607557$  & $-0.1657057$ \\
80  & $-8508.29536397$  & $-1.3294212$ & $7674.96382916$  & $-0.1302081$ \\
120 & $-19408.86914398$ & $-1.3478381$ & $17604.86199438$ & $-0.1252783$ \\
160 & $-34743.05318075$ & $-1.3571505$ & & \\
200 & $-54510.86712322$ & $-1.3627717$ & & \\
240 & $-78712.313$\phantom{00000} & $-1.36653$\phantom{00} & & \\
\bottomrule
\end{tabular}
\caption{Exact determinant data. $G_K(\z)$ is positive definite at all six $K$
(Proposition 2.2), and (6.16) holds at all six. $P_K=\Delta_K/\operatorname{cont}(\Delta_K)$ is
the primitive polynomial, so $S_K$ cancels identically from $K^{-2}\log P_K(\z)$, and the
fourth column is $\log(1/\operatorname{cont}(F_K))$, i.e.\ the true minimal normalizer.
$K\le120$ was produced by four independent implementations and $K\le200$ by two, each at two
precisions; the $K=240$ row comes from a single track at $400$--$800$ digits and is quoted to
its stated precision.}
\label{tab:det}
\end{table}

The last column of Table~\ref{tab:det} is, in our view, \textbf{the single most informative
number in the audit}. It is the decay actually achieved by the best possible normalization, and
it must be compared with the requirement $<-139/8000=-0.0173750$. The measured values
$-0.166,-0.130,-0.125$ extrapolate to roughly $-0.12$ (plausible range $[-0.27,-0.072]$), and
the value predicted from the verified asymptotics is
$A_\infty+\Ub_{\mathrm{true}}=1.3134540-1.3855098=-0.0720558$. \emph{The construction has about
seven times the room it needs.} For calibration, the constructions in this circle of ideas that
are known to fail do so by a positive margin: the Brown--Zudilin $\zeta(5)$ cellular forms give
$+0.22$ per $N^2/2$ with true denominators, and Rivoal's $(3;1^q)$ families at $q=11$ miss by
$+0.05$ to $+0.14$. This construction is not marginal in the way those are; only the paper's
\emph{printed exponent} is marginal.

\subsection{The local bounds against the truth}

\begin{table}[htbp]
\centering
\begin{tabular}{@{}rl@{}}
\toprule
$K$ & $\vG(\Delta_K)-\gout$ at the primes satisfying (4.9), in increasing order of $p$ \\
\midrule
40  & $1,1,0,0,0,0$ \hfill ($p=17,\dots,37$) \\
80  & $39,15,1,0,0,0,0,0,0,0,0,0,0$ \hfill ($p=29,\dots,79$) \\
120 & $78,60,17,1,1,0,0,0,0,0,0,0,0,0,0,0,0,0$ \hfill ($p=41,\dots,113$) \\
\bottomrule
\end{tabular}
\caption{Proposition 4.3 against exact determinants: 0 violations at 37 primes, equality at 27.
The slack that does occur is confined to $K\ge2p$ and is in the safe direction.}
\label{tab:gout}
\end{table}

For the inner range the corresponding test is impossible, and Table~\ref{tab:inner} records what
was done instead.

\begin{table}[htbp]
\centering
\begin{tabular}{@{}rrrr@{}}
\toprule
$L_0$ & side conditions & entry-bound failures & $\gin$ (truth: $-188$) \\
\midrule
0 & zero-source dominance fails & 2094 & $-170$ \quad(\emph{false bound}) \\
1 & fails & 812 & $-180$ \quad(\emph{false bound}) \\
2 & fails & 29 & $-186$ \quad(\emph{false bound}) \\
3 & all hold & \textbf{0} & $\mathbf{-188}$ \quad(exact) \\
4 & all hold & \textbf{0} & $\mathbf{-188}$ \quad(exact) \\
\bottomrule
\end{tabular}
\caption{The mechanism of Proposition 4.1 at the smallest nondegenerate instance, $K=80$,
$p=23$, where $v_{23}^{\mathrm{G}}(\Delta_{80})=-188$. The paper's side conditions are exactly
what separates a true bound from a false one.}
\label{tab:inner}
\end{table}

Finally, the limiting functions really are the limits: $\gin-p\Gamma(K/p)$ is bounded by $3.8M$
uniformly in $K$ and is \emph{exactly} 35 at fixed $x\approx10$, $M=40$, for
$K=3.2\times10^5,\dots,1.024\times10^7$; $v_p(S_K)-p\mathcal{N}(K/p)$ is likewise $K$-independent
at fixed $x$; and $-\gout-K(R_0-d)(p/K)$ has modulus at most 7. Hence $A_M$ is the right
arithmetic constant for the local bounds the paper proves.

\section{What would have to be true for the proof to work}\label{sec:musthold}

Stripped of everything that has been verified, the paper rests on the following load-bearing
statements. This is the list the author should be asked to defend.

\begin{enumerate}
\item \textbf{$\vG(\Delta_K)\ge\gin$ for the inner primes $K/M<p\le K/3$ (Proposition 4.1).}
The one substantive claim no computation in this audit could reach. Everything around it checks
out --- the combinatorics at the real parameters, the mechanism at $K=80$, $p=23$, the
hand-off degrees, the limit $\Gamma$ --- but the inequality itself is unverified at admissible
parameters, and if $\gin$ over-claims, then $Q_{K,M}\notin\mathbb{Z}[X]$ and both Theorem 1.1
and Corollary 1.2 collapse silently.

\item \textbf{Lemmas 3.1 and 3.2, at the parameters \S4 uses them.} A one-power-of-$p$ error
would move $\gin$ by about $h$ and change $A_M$. Both lemmas were verified exactly at small $p$;
the necessity of $d(r_\nu)<p$ and of $p\ge7$ was confirmed by counterexample.

\item \textbf{$b_a\le6\alpha x+3$ (p.~10).} True and sharp, but unproved in the paper and
load-bearing for $L_a\ge0$ at $x=3$; see Finding~\ref{f:ba}. Two lines would fix it.

\item \textbf{The $m_N\ge1$ part of the inner range.} For $M=200$ this is most of the inner
range and it has no determinant test at any $K$. The bookkeeping is uniform in $m_N$ and was
verified symbolically and over sweeps reaching $m_N=14$.

\item \textbf{$A_{200}+\Ub<0$.} Verified exactly, \emph{given} $\gin$ and $\gout$. The margin is
$1.29\%$; an error of that size in either constant kills Theorem 1.1.
\end{enumerate}

\noindent\textbf{A structural remark that should reassure the author and the reader alike.} By
the paper's own C.3, $Q_{K,M}=a_{K,M}P_K$ with $a_{K,M}$ a positive integer and
$P_K=\Delta_K/\operatorname{cont}(\Delta_K)$ primitive. Hence Theorem 1.1 follows from
\[
\limsup_{K\to\infty,\;40\mid K}K^{-2}\log P_K(\z)<0
\]
alone, with no reference to $m_{K,M}$, to $\gin$, or to Sections 3--5 at all: $P_K\in\mathbb{Z}[X]$
has degree $h$ by (2.9), is positive at $\z$ by Proposition 2.2, and the same final step applies.
The entire arithmetic apparatus exists only to \emph{prove} such a bound, by exhibiting an
explicit divisor $m_{K,M}$ of the content. The measured values in Table~\ref{tab:det} say that
the quantity above is about $-0.12$ at computable $K$ and extrapolates to $-0.072$. So an error
in Sections 3--5 would be an error in the proof and quite probably not in the theorem. We say
this not to excuse a gap but to locate it: a referee's report on this paper should be about
whether the denominators are proved, not about whether the construction works.

\section{Minor and presentational remarks}

These are offered to the author as corrections; none affects the mathematics.

\begin{enumerate}
\item \emph{(\S2.)} ``Polynomial division and simple partial fractions extend $\mu_X$'' (p.~3) is
a definition on a basis, not an extension; the exclusion of $t=0$ from the pole set is used ---
it is what makes Proposition 2.2's integral converge --- but never stated. In Proposition 2.2
(p.~5) the interchange of $\sum_{\ell\ge1}$ and $\int_0^\infty$ is not justified (Tonelli, one
line). In \S2.3 (p.~4) the hedge ``the \emph{possible} poles $-j^2$'' is confusing, since all $h$
residues are nonzero, which is what makes $V$ square and (2.9) nonvanishing; and $W=D_N^5$ is
defined there but first used in \S4.
\item \emph{(\S3.)} On p.~6 the citation ``[9, (24.4.1), (24.4.3)]'' appears to list the two
DLMF displays in the opposite order to the two identities named in the sentence. We were not
able to confirm the DLMF numbering (no offline copy; see Section~\ref{sec:notchecked}), and in
any case a bracketed list of two references carries no ordering, so this is at most a
readability point. Lemma 3.3's residue sign is left as
``$\pm$''; it is $(-1)^{K-r}$ uniformly. In Lemma 3.1 the hypothesis $\deg U_0\le p+1$ may be
relaxed to $4p-2$ and $r_\nu-r_\eta\in\mathbb{Z}_p^\times$ is redundant --- unused slack, worth
knowing if \S\S4--5 ever needs it.
\item \emph{(pp.~10, (4.2)--(4.3).)} The ``$-4$'' and ``$+1$'' silently absorb the $p^{-4}$
prefactor of (3.7); the text does not say so, and settling it costs the reader a full
re-derivation.
\item \emph{(p.~12.)} The divided-difference integrality needs the $-\tfrac14+\tfrac1{2j}$ tail
of (2.3), via $a^4(H^{(5)}_{a-1}-H^{(5)}_a)+1/a=0$. The paper points at (2.3), so the reader has
what is needed, but the one-line computation is missing; over 2087 node pairs the congruence
holds $2087/2087$ with the tail and $0/2087$ without.
\item \emph{(p.~14, (5.7).)} No explicit constant is given for the $O_M(1)$; the correct one is
$3.8M$. ``There are $O(M)$ such intervals'' over-counts: in the $z$ variable there are at most
three pieces. The cancellation between the zero block and the ordinary block should be described
as a cancellation rather than as an omission.
\item \emph{(p.~23, (A.6).)} ``The parenthesis is negative'' follows in one line from the
preceding sentence $\Phi'(\sqrt{q_-})<0$ and the monotonicity of $\arctan$; as printed it looks
like an unsupported assertion, and it holds by only $2.5\times10^{-8}$.
\item \emph{(pp.~24--25, Table 2.)} The caption explains only the single-run format. Rows with
two runs are momentarily unreadable; one clause --- that runs are separated by a space --- would
fix it.
\item \emph{(p.~26, (B.2).)} Redundant: $2(1-\alpha)=2\lambda$ appears twice and the
$2\alpha$-points are a subset of the $4\alpha$-points. Also, $H=\lambda+3\alpha$ should be stated,
since it is what makes the quadratic coefficient of $R$ vanish identically.
\item \emph{(Appendix C.)} The displayed $h^2\mathcal{R}^{2h}$ in (C.5) is true and reproducible
but weaker by the $O(1)$ factor $\mathcal{R}^2=293.84$ than the same step gives. ``Its positive
margin absorbs the polynomial factor in $h$'' (p.~30) needs $K\gtrsim2.81\times10^5$; neither the
margin $1.433\times10^{-4}$ nor the threshold is stated. And ``the threshold $b_0$ is not
specified'' understates: $b_0$ is not \emph{effective}, and is at least $\exp(1.07\times10^{11})$.
\item \emph{Notation.} $s$ is the number of subintervals in B.1, the quantity $Hx-T/2$ in (5.4),
and appears in $O(K^{3/2})$; $d$ is $d(r)$ in \S3, $d(y)$ in (5.9) and $d_0$ in B.1; $M$ is the
cutoff while $M_0$ is a potential constant; $N$ is $3n$ while $\mathcal{N}(x)$ is (5.5).
\item \emph{Robustness (the most valuable single change).} Quoting (6.2) with
$M_0=-6.6498$ instead of $-6.645$ --- which the author's own data support --- would multiply the
final margin in (7.2) by $139$ (from $0.0525$ to $7.296$) and remove the knife-edge from (6.4),
(A.9) and the printed exponent at a stroke. Similarly, the sanity check that the machinery gives $+0.93$ for $\zeta(7)$
and $-0.86$ for $\zeta(3)$ is half a page and answers the first question every referee will ask.
\item \emph{``Author's note on further work'' (p.~31).} A page of unsubstantiated claims of
``several advances'' and ``substantial progress'' with no mathematical content, which the section
itself then disclaims. It raises a referee's prior against the paper for no gain and should be
deleted.
\item \emph{Reproducibility.} The paper should ship an ancillary verification script. Everything
in Appendices A and B is reproducible at ordinary precision in minutes --- we did it --- but a
referee should not have to reimplement the rational-enclosure apparatus to check a lemma.
\end{enumerate}

\section{What we did not check, and the limits of this audit}\label{sec:notchecked}

\subsection{Not checked}

\begin{enumerate}
\item \textbf{$\vG(\Delta_K)\ge\gin$ at admissible parameters.} The central omission. Hypothesis
(4.1) forces $h\ge296000$; the smallest admissible matrix is $296000\times296000$. What would
settle it: a machine-checked formalisation of \S3 plus \S4.1 (the argument is finite and
elementary), or a targeted determinant over $\mathbb{Z}_p$ in the unimodular basis (4.5) at a
single admissible prime, or --- much the cheapest --- a rewrite of \S4 admitting small $M$
(e.g.\ $M=3$, $K\approx1800$) so that the prime-by-prime comparison becomes computable.
\item \textbf{The $m_N\ge1$ portion of the inner range}, verified only by symbolic bookkeeping.
\item \textbf{The author's own exact-rational enclosure pipeline} over the 684 intervals of
Table 2. We verified the same inequalities in 25--60-digit interval-free arithmetic against a
smallest margin of $6.9\times10^{-7}$, and verified analytically that the prescribed truncation
orders leave remainders $10^{37}$ times smaller than needed. The residual risk is negligible but
not zero; an interval-arithmetic or formally verified rerun would close it.
\item \textbf{An independent transcription of Tables 1 and 2.} Ours checksums (exact mass
$37/40$; exact tiling of $[0,2]$), which is strong evidence, but a second pair of eyes on
Table 1 is worth one referee-hour given the $6.45\times10^{-8}$ margin in (6.4).
\item \textbf{Every proof step that is not a computation.} The audit is a machine audit with a
line-by-line reading on top of it. The proofs of \textbf{Lemmas 3.1, 3.2, 3.3}, of
\textbf{Lemma 4.2}, of \textbf{Lemma 6.2}, and the partial-summation step inside
\textbf{Proposition 5.2} were re-derived by hand and then tested only through their
\emph{consequences} (exactly, at small $p$ and small $K$, and by randomised stress tests). None
of them has been verified as a theorem at the parameters \S\S4--5 use them. We found no defect
in any of them, and Lemma 4.2's statement is a finite claim about $h\times h$ matrices that
6000 random instances did not break; but a machine-checked proof is the thing that would settle
them.
\item \textbf{The degree bound on p.~9} (``the pullback numerator has degree at most
$12N+4h+1$''), an input to Lemma 3.3 and hence to (3.12). We confirmed the \emph{use} made of
it --- $\max(2K,d+1)=184n+2\le200n=5K$, and replacing $\max(2K,d+1)$ by $5K$ weakens (3.10),
which is the safe direction --- but not the bound itself.
\item \textbf{The DLMF equation numbers} cited on pp.~5, 6 and 29 ((25.6.2), (25.11.29),
(24.4.1), (24.4.3), \S24.10(i), (18.10.5)). We verified the \emph{identities} numerically
(Euler's $\zeta(2m)$ formula symbolically; Hermite's formula to $\ge24$ digits; von
Staudt--Clausen through $v_p(\kappa_d)\ge-1$; the Legendre bound $|P_j(s)|\le R^j$ for
$j\le24$); we did not open DLMF to confirm that those are their numbers.
\item \textbf{The $(\alpha,c)$ optimisation behind the $\zeta(s)$ sanity test.} The table of
balances at the paper's own $(\alpha,\lambda)=(3/40,37/40)$ was reproduced exactly by two
tracks; the further claim that re-optimising $(\alpha,c)$ still leaves $\zeta(7)$ failing
($+0.84$) rests on one track's computation and was not independently re-run.
\item \textbf{Whether a known no-go result applies.} Zudilin's $\mathbb{Z}$-transfinite-diameter
obstruction, F.~Brown's recent bound and the Chudnovsky--Bombieri $G$-function bounds all
constrain linear-form or Mellin constructions; none that we found applies verbatim to a
determinant producing integer \emph{polynomials} $Q$ with $Q(\z)$ small. A short answer from a
specialist (Zudilin, Brown, Fischler) is the one external input that could settle the paper
quickly in either direction. Separately, the author and this construction appear nowhere we
could find. The bibliography [1]--[13] is a list of real papers; we checked four of them at
source ([4] Brown--Zudilin, [5] Fischler, [6] Fischler--Sprang--Zudilin, [8] Lai--Lupu--Sprang)
and the paper characterises each accurately, including the two quotations on p.~1.
\end{enumerate}

\subsection{Where the audit corrected itself}\label{sec:disagree}

Discipline requires recording where our own first pass was wrong, since several of these would
have become false accusations against the paper.

\begin{itemize}
\item \textbf{Table 3, row $j=14$ (p.~27).} One verification pass reported that the paper prints
$167757646109/955086612480$ where the truth is $16775764609/955086612480$, and called it ``the
only outright false printed statement in the paper''. \textbf{This is wrong, and we settled it
ourselves.} The PDF text layer reads $16775764609$, and so does a 600\,dpi render of p.~27,
where the numerator is unambiguously eleven digits long; it is only at screen resolution that
the glyphs invited a misread, which two separate passes then made in the same way. The decisive
test is arithmetic rather than typographic: only $16775764609/955086612480$ makes the 17 printed
rows sum to (5.18), and with it the sum is \emph{exactly}
$322437603634266857629/7535670527041937280000$ (recomputed here in \texttt{Fraction}
arithmetic). \textbf{The paper is correct here, and its sentence ``Grouping its terms into unit
intervals gives Table 3. Their sum is (5.18)'' is true as printed.} General rule for anyone
re-doing this work: take digits from the text layer or from a high-resolution render, never from
a screen-resolution page image.
\item \textbf{The (A.9) margin.} One track reported $6.892\times10^{-10}$; its verifier reported
$6.892\times10^{-7}$. We recomputed it from our own parse of Table 2: the margin is
$6.8921559\times10^{-7}$. The verifier is right.
\item \textbf{``Table 2 is unparseable, so (A.9) cannot be checked''} (recorded as a major
finding on the first pass). Refuted: the $k$-column lists whitespace-separated runs, the 684
intervals tile $[0,2]$ exactly, and (A.9) holds on all of them. Lemma 6.1 moves from ``could not
check'' to \emph{verified}. Only the caption needs a clause.
\item \textbf{``The Bernoulli convention $B_1=-1/2$ is never stated.''} False: it is stated on
p.~3, one line before (2.2). We confirmed this in the text layer.
\item \textbf{``$p^2>5K$ is asserted but never used.''} False: it is used on p.~14 in \S5.1, to
truncate Legendre's formula (5.3) to its $a=1$ terms.
\item \textbf{``Proposition 4.1 relies on three side conditions never displayed.''} False: all
three are displayed in the proof on pp.~10--11. They are consequences of (4.1) proved inside the
proof, not hidden assumptions. (The genuine gap in that stretch is the separate
$b_a\le6\alpha x+3$, Finding~\ref{f:ba}.)
\item \textbf{``pdftotext garbles (4.9) as $5K\le2p-2$.''} False for our text extraction:
line 613 of it reads $5N\le2p-2$, matching the PDF. Publishing this caution would have sent
readers chasing a non-existent problem.
\item \textbf{``(3.8)'s proof suppresses a case split.''} A misreading: the paper proves (3.8) by
Vandermonde together with $\binom{z}{j}=\tfrac{z}{j}\binom{z-1}{j-1}$, which needs no case split.
\item \textbf{The uniformity constant in (5.7).} One track reported
$0<\gin-p\Gamma<7.41M^2$, always positive; its verifier traced this to a factor-2 bug in that
track's code (a cap on the zero-block weights applied to an already-doubled quantity) and
obtained $|\gin-p\Gamma|\le3.8M$ with both signs. The verifier reproduced the erroneous numbers
with the buggy line and the correct ones without it. The paper is unaffected, and is in fact
better served by the corrected statement.
\item \textbf{Smaller corrections} adopted from verification passes: the sharp constant in (6.11)
is $(2\pi)^4/12=129.88$, not $123.575$ (a grid artefact); $C_p\in p^5\mathbb{Z}_p^\times$ is false
at $p=37$, the correct statement being $v_p(C_p)\ge5$ with equality unless $p\mid B_{p-5}$;
Proposition 4.3's bound is attained at 27 of 37 primes, not 26; $u>0$ occurs already at $K=80$;
and (C.5)'s finding was downgraded to cosmetic once it was seen that the paper's displayed form
\emph{is} reproducible from its own argument.
\end{itemize}

\subsection{Limits}

This audit is numerical and structural, not a formal verification. Its strongest claims are of
the form ``this identity is exact'' and ``this formula is exactly attained at 27 independent
primes''; its weakest concern the asymptotic regime the paper actually inhabits
($K\ge8\times10^6$ for $M=200$, $K\ge2\times10^{12}$ for Corollary 1.2). Every extrapolation we
report --- $\Ub_{\mathrm{true}}$, the decay of $P_K(\z)$, the constant $A_\infty$ --- is a fit
over $K\le240$, validated out of sample where possible but not proved; and where the paper's
statements are $\limsup$s with no effective threshold, no finite computation can confirm them, so
we checked the constants they are made of instead. What the audit did \emph{not} consist of is
reading the paper for plausibility: every track reproduced a known answer --- Euler's
$\zeta(2m)$, Hilbert determinants, Marchenko--Pastur, von Staudt--Clausen, Raabe, DLMF 25.11.29
--- before it was allowed a verdict, and was then attacked by a second independent
implementation. In every case, the paper survived.

\section*{Provenance and authorship}
\addcontentsline{toc}{section}{Provenance and authorship}

This report was written by \textbf{Claude Opus 5 (1M context)} (model identifier
\texttt{claude-opus-5[1m]}), an AI system made by Anthropic, running in Claude Code on Dan
Romik's machine on 22 September 2026. No part of it has been refereed by a human mathematician.
It is signed so that a reader can weigh it accordingly, and so that any error in it is
attributable.

\medskip
\noindent\textbf{How it was produced.} The audit ran as an orchestrated workflow of twenty agent
instances of the same model: nine audit tracks --- the construction and Appendix~C; the local
functional (\S3); the local determinant estimates (\S4); normalization and the prime sum
(\S5, Appendix~B); the real bound (\S6, Appendix~A); exact computation of $\Delta_K(X)$; independent
recomputation of the two constants; a first-principles equilibrium analysis; and the strategy,
logic and literature --- each followed by a \emph{separate} adversarial verification pass
instructed to attack the track's findings from both sides, re-implementing the decisive
computations by a different method; then a writer and a fact-checker. Each track was required to
reproduce a known answer before it was permitted a verdict. The verification passes materially
changed the report: they withdrew two of the \S4 findings as mistaken, moved Lemma 6.1 from
``could not check'' to verified, and refuted a claimed transcription error in Table 3 by
rendering the page at 600\,dpi. The fact-checking pass found twenty further defects in the draft,
five of them substantive, including two wrong final digits in $A_{200}$ and $A_{100000}$ and a
confidence figure higher than the weakest component it aggregated.

\end{document}
